% Chapter 1
% 
\chapter{Introduction} % Main chapter title
\label{chap:introduction} % For referencing the chapter elsewhere, use Chapter~\ref{chap:introduction}

The problem addressed by this thesis and the reasoning behind it are introduced in this chapter. Section~\ref{sec:context_motivation} describes the context and motivation, Section~\ref{sec:hypothesis} states the research hypotheses, and Section~\ref{sec:questions} lists the research questions. The general and specific objectives follow in Section~\ref{sec:objectives}, and Section~\ref{sec:structure} closes the chapter with an overview of how the rest of the document is organized.

%-------------------------------------------------------------------------------
%---------
%
\section{Context and Motivation} 
\label{sec:context_motivation}

Identifying specific errors and learning difficulties in student work is a time-consuming and challenging task for educators across primary and secondary education (K--12). Teachers must routinely interpret short answers, essays, worksheets, and problem-solving traces while balancing instruction, planning, and administrative work. Large-scale workload evidence shows why timely identification is difficult in practice: in the OECD Teaching and Learning International Survey (TALIS) 2018, lower-secondary teachers report spending a substantial working week on non-teaching tasks, including several hours per week on marking and correcting student work \parencite{OECDTALIS2018VolI}. The 2024 round of the survey extends this evidence on the demands of teaching and also reports whether and how teachers use artificial intelligence \parencite{OECDTALIS2024}. In parallel, class size in lower secondary education commonly exceeds 20 students and is considerably larger in some education systems, which increases the volume and heterogeneity of student responses that need assessment \parencite{OECDEAG2019}. When assessment occurs manually at scale, feedback tends to focus on surface correctness and completion; recurring error mechanisms that indicate persistent difficulties receive less attention. As a consequence, patterns that inform targeted support often become visible only after repeated failures or summative assessments, despite long-standing evidence that earlier identification and support prevent compounding learning gaps \parencite{SnowBurnsGriffin1998}.

A non-trivial fraction of learners experience persistent difficulties in foundational skills. Dyslexia (a specific reading disability), for example, is far from rare. A common definition, based on a reading-achievement cutoff of 1.5 standard deviations below the mean for age, identifies about 7\% of the population as dyslexic, and estimates vary with the diagnostic criteria adopted \parencite{PetersonPennington2012}. At the system level, international learning indicators point to the urgency of scalable analytical support. The World Bank has reported high rates of ``learning poverty'' (the inability to read and understand a simple text by age 10) in low- and middle-income contexts, which shows how widespread foundational difficulties become when early support is constrained \parencite{WorldBank2022LearningPoverty}. These figures motivate tools that help educators detect patterns earlier and more consistently, without replacing professional judgment or producing clinical diagnoses.

Artificial Intelligence (AI) and, more specifically, Natural Language Processing (NLP) and learning analytics offer a concrete path to scale formative assessment. In writing tasks, this goal has driven Automated Writing Evaluation (AWE) systems and error detection research that combine scoring with feedback-oriented language analysis \parencite{Bryant2023GECsurvey}. Research on error detection and Grammatical Error Correction (GEC) has benefited from shared tasks and benchmarks that enabled comparable evaluation \parencite{Ng2014,Bryant2019}. Yet these resources and the strongest-performing models are unevenly distributed across languages and educational domains, which complicates generalizable deployment in diverse educational settings \parencite{Joshi2020}. For Portuguese, recent work includes an LLM-assisted characterization of errors in a European Portuguese learner corpus \parencite{Akef2026} and a synthetic dataset for grammar correction in Brazilian Portuguese \parencite{Cabral2026}. In educational workflows, correcting a single response is rarely enough; teachers need summaries that connect error types to likely sources of difficulty and that support monitoring progress over time \parencite{RomeroVentura2024}. Large Language Models (LLMs) can generate controlled synthetic data for rare error categories and help scale development, but the educational use case demands explicit constraints on reliability and transparency, especially when the output informs instruction, not grading \parencite{Brown2020}. Recent work on the use of LLMs in education reviews both the opportunities they open for learning and teaching and the challenges that accompany them \parencite{Kasneci2023}.

In mathematics and other structured-response tasks, formative assessment goals commonly go beyond correctness to identifying systematic misconceptions and recurring procedural ``bugs'' that guide targeted remediation \parencite{VanLehn1990}. The opportunities for scalable text-based analysis are especially relevant in digital learning environments, where students submit typed responses (short answers, explanations, essays, and numeric solutions) through learning management systems, online assessments, and digital worksheets. This thesis focuses exclusively on such digital, text-only workflows. Authentic classroom assessments often include handwritten responses, but integrating handwriting recognition (OCR/HTR) introduces distinct technical challenges that are outside the scope of this work and are reserved for future research. By isolating the text analysis component, this thesis aims to establish a baseline for error detection in digital environments.\par

The main gap addressed in this thesis is the lack of teacher-oriented systems that (i) detect and classify error patterns specifically linked to learning difficulties (e.g., phonological, orthographic, and segmentation errors in language; procedural bugs and conceptual misconceptions in mathematics), (ii) aggregate these errors over time into interpretable difficulty profiles, and (iii) do so under realistic data-scarcity constraints using synthetic data augmentation. The intended outcome is actionable, non-clinical information that supports educators' formative assessment and planning without providing clinical diagnoses of learning disabilities.

\section{Research Hypothesis}
\label{sec:hypothesis}

Based on these challenges, the following hypotheses are proposed. Throughout this thesis, performance is assessed using standard evaluation measures from educational data mining (EDM) and NLP (precision, recall, \(F_1\)-score), complemented by clustering validity indicators. Inter-rater agreement measures such as Cohen's \(\kappa\) \parencite{Cohen1960} would apply to a teacher-annotation study, and qualitative educator review to a classroom pilot; neither was conducted within this thesis, and both are specified as future work (Chapter~\ref{chap:Chapter6}). The hypotheses are intentionally formulated at a high level so that model choices, integration strategies, and evaluation protocols evolve as empirical evidence and design constraints emerge during the thesis.

\begin{itemize}
    \item \textbf{Main Hypothesis (H1):} Text-based NLP models and symbolic methods, trained on a combination of authentic and synthetic student responses, can identify pedagogically meaningful error types (phonological, orthographic, and segmentation errors in language; procedural bugs and conceptual misconceptions in mathematics) with performance consistent with expert teacher annotations and established instructional taxonomies, without providing clinical diagnosis. Performance is evaluated on typed inputs using standard classification metrics (precision, recall, \(F_1\)-score).\footnote{No new teacher-annotation study was conducted; the evidence in Chapter~\ref{chap:Chapter5} measures performance against the expert-produced gold annotations of existing corpora, which tests the detection component of H1 but not its consistency with independently collected teacher judgments.}
    \item \textbf{Secondary Hypothesis (H2):} Using rule-based generators and Large Language Models (LLMs) to produce synthetic learner-error data improves the robustness and generalization of error detection and classification models in data-scarce educational scenarios. The hypothesis thus seeks to validate, for current students and in the present educational context, the benefits reported for synthetic training data in the literature (Section~\ref{sec:sota_synth}). Improvements are assessed empirically using held-out evaluation and generalization checks (e.g., cross-validation and transfer across tasks or difficulty levels).\footnote{As documented in Section~\ref{sec:method_data_llm}, LLM-based generation was considered and deliberately not adopted; the empirical evidence in Chapter~\ref{chap:Chapter5} therefore addresses the rule-based part of this hypothesis and leaves the LLM part untested.}
    \item \textbf{Secondary Hypothesis (H3):} Unsupervised pattern discovery methods aggregate recurring errors over time into meaningful student difficulty profiles that are interpretable and actionable for educators. Here too, the aim is to validate, for current students and in the present educational context, the profiling approaches reported in the EDM literature (Section~\ref{sec:sota_ethics}). Profile quality is assessed through a combination of clustering validity indices and pedagogical interpretability.\footnote{Interpretability is assessed by inspection of cluster centroids against the taxonomy (Section~\ref{sec:method_models_profile}); validation of pedagogical meaningfulness and actionability with educators was not conducted and is part of the pilot proposed in Chapter~\ref{chap:Chapter6}.}
\end{itemize}

\section{Research Questions}
\label{sec:questions}

To test the proposed hypotheses and guide the development of the system, this research seeks to answer the following questions:

\begin{enumerate}
    \item[\textbf{Q1}] Linguistic error taxonomies: What linguistic error types (phonological, orthographic, and segmentation) are observable and recurrent in the available typed learner corpora, how do their distributions differ across learner populations, and how are they best structured into a formal taxonomy suitable for automated text analysis?
    \item[\textbf{Q2}] Detection and classification models: To what extent do NLP sequence-labeling models automatically detect and classify difficulty-linked linguistic errors in typed text, as measured by span-level \(F_1\)-score against expert-produced gold annotations, including on text from the target population?
    \item[\textbf{Q3}] Mathematical misconception detection: How well do symbolic methods and simple ML classifiers identify characteristic arithmetic procedural bugs (e.g., smaller-from-larger subtraction) and fraction misconceptions (e.g., whole-number bias) from typed numeric answers?
    \item[\textbf{Q4}] Data scarcity and synthesis: In the context of limited annotated educational datasets, what combination of real student data, rule-based synthetic error generation, and LLM-based augmentation provides the most effective training set, as measured by \(F_1\)-score stability and generalization, while maintaining fidelity to authentic learner error distributions?
    \item[\textbf{Q5}] Difficulty profiling: Do aggregated error patterns cluster into stable and interpretable student difficulty profiles, and how might these profiles be presented to educators in a non-clinical format?
\end{enumerate}

\section{Objectives}
\label{sec:objectives}

Concrete goals follow from the hypotheses and questions above. A general objective states the overall aim of the work, and six specific objectives (O1--O6) split it into the tasks addressed in the following chapters, covering the error taxonomies, the data strategy, model development, and evaluation.

\subsection{General Objective}
The primary objective of this thesis is to design and evaluate an NLP-based system capable of analyzing typed student responses in primary and secondary education (K--12) to detect specific error patterns linked to learning difficulties and to infer interpretable difficulty profiles. The system should provide educators with actionable, non-clinical information to support formative assessment and pedagogical planning.

\subsection{Specific Objectives}
\begin{enumerate}
    \item[\textbf{O1}] Define and formalize error taxonomies. Identify and categorize linguistic errors (phonological, orthographic, and segmentation) and mathematical errors (procedural bugs in arithmetic, fraction misconceptions) into formal taxonomies suitable for automated detection.
    \item[\textbf{O2}] Develop a data strategy combining real and synthetic data. Design rule-based generators to produce synthetic student responses with controlled error patterns, complementing limited real educational data (LLM-based generation was considered and not adopted; Section~\ref{sec:method_data_llm}).
    \item[\textbf{O3}] Develop linguistic error detection models. Implement and evaluate NLP sequence-labeling models (e.g., fine-tuned Transformer encoders) for token-level or span-level error tagging on typed text.
    \item[\textbf{O4}] Develop mathematical misconception detection. Implement symbolic computation and template-matching methods to identify procedural bugs and conceptual misconceptions from typed numeric answers.
    \item[\textbf{O5}] Implement student difficulty profiling. Aggregate error statistics per student and apply clustering algorithms to discover interpretable difficulty profiles.
    \item[\textbf{O6}] Evaluate system performance and pedagogical utility. Define metrics (precision, recall, \(F_1\), clustering indices) for quantitative evaluation and assess interpretability for teacher-facing use by inspection of the resulting profiles (educator validation is deferred to a school pilot).
\end{enumerate}

\section{Thesis Structure}
\label{sec:structure}

The remainder of this dissertation is organized as follows. Chapter~\ref{chap:Chapter2} presents the state of the art in NLP for education, grammatical error correction and detection, mathematical error analysis, synthetic data generation for educational AI, and the pedagogical and ethical dimensions of learning-difficulty identification. This chapter positions the thesis with respect to the main gaps: (i) the underrepresentation of certain languages and local educational contexts; (ii) the limited use of synthetic data for difficulty-linked error patterns; and (iii) the tendency to focus on isolated responses, with longitudinal difficulty patterns receiving less attention.

Chapter~3 turns to the proposed error taxonomies and the data strategy. It formalizes linguistic error categories (phonological, orthographic, segmentation) and mathematical error categories (procedural bugs, fraction misconceptions). The data approach combines limited real student data with rule-based synthetic error generation (LLM-based generation was considered and deliberately not adopted), and it includes quality control procedures.

Chapter~4 covers the system architecture and its implementation: the NLP-based linguistic error detection module, the symbolic/template-based mathematical misconception detection module, and the profiling module that aggregates errors into student difficulty profiles. Implementation details are documented there.

Chapter~5 contains the experimental evaluation. Quantitative results cover error detection (precision, recall, \(F_1\), with paired bootstrap confidence intervals and multi-seed replication), domain shift to the target population, the impact and scaling behavior of synthetic data (including a real-only baseline trained without any synthetic stage), mathematical misconception detection under clean and perturbed conditions, and difficulty-profile validity (clustering recoverability and split-half stability via the Adjusted Rand Index).

Chapter~6 concludes the thesis by summarizing contributions, discussing limitations, and outlining future work, including extension to additional languages and real classroom deployment. The bibliography and appendices follow.